\documentclass[a4paper, 10pt]{article}
\usepackage{scrextend}
\changefontsizes{9pt}
\input{../course-config}
\usepackage{../vendor/template/CEDT-Assignment-style}

\begin{document}
\subject
\asgntitle{6}{}{Week 6}{\StudentID\ \StudentName}{\DefaultCollaborators}

% ================================================================================ %
%                                    Problem 01                                    %
% ================================================================================ %
\begin{problem}
Design a Self-Checking 2:4 Demultiplexer with Enable.

\begin{figure}[h]
	\centering

	\begin{minipage}[c]{0.54\textwidth}
		\centering
		\resizebox{\linewidth}{!}{%
			\begin{tikzpicture}[
					circuit logic US,
					wire/.style={
							draw,
							line width=0.8pt,
							line cap=round,
							line join=round
						},
					junction/.style={
							circle,
							fill=black,
							inner sep=1.5pt
						},
					gate/.style={
							draw,
							line width=0.8pt,
							minimum width=1.25cm,
							minimum height=0.95cm,
							logic gate inputs=nnn
						}
				]

				% =========================================================
				% AND gates
				% =========================================================
				\node[and gate US, gate] (g0) at (7.6,6.0) {};
				\node[and gate US, gate] (g1) at (7.6,4.0) {};
				\node[and gate US, gate] (g2) at (7.6,2.0) {};
				\node[and gate US, gate] (g3) at (7.6,0.0) {};

				% =========================================================
				% Reference levels
				% =========================================================
				\coordinate (Alevel)  at (0,0 |- g0.input 1);
				\coordinate (Blevel)  at (0,3.0);
				\coordinate (ENlevel) at (0,0 |- g3.input 3);

				% =========================================================
				% Input labels
				% =========================================================
				\node[anchor=east] (Alabel)  at (-0.15,0 |- Alevel)  {$A$};
				\node[anchor=east] (Blabel)  at (-0.15,0 |- Blevel)  {$B$};
				\node[anchor=east] (ENlabel) at (-0.15,0 |- ENlevel) {$EN$};

				% =========================================================
				% NOT gates
				% =========================================================
				\node[
					not gate US,
					draw,
					line width=0.8pt,
					minimum width=0.75cm,
					minimum height=0.55cm
				] (notA) at (1.55,0 |- Alevel) {};

				\node[
					not gate US,
					draw,
					line width=0.8pt,
					minimum width=0.75cm,
					minimum height=0.55cm
				] (notB) at (1.55,0 |- Blevel) {};

				% =========================================================
				% A input
				% =========================================================
				\coordinate (Asplit) at (0.75,0 |- Alevel);

				\draw[wire]
				(Alabel.east)
				-- (notA.input);

				\node[junction] at (Asplit) {};

				% =========================================================
				% A' distribution
				% =========================================================
				\coordinate (AbarTop)    at (3.05,0 |- g0.input 1);
				\coordinate (AbarBottom) at (3.05,0 |- g2.input 1);

				\draw[wire]
				(notA.output)
				-- (g0.input 1);

				\draw[wire]
				(AbarTop)
				-- (AbarBottom)
				-- (g2.input 1);

				\node[junction] at (AbarTop) {};

				% =========================================================
				% A distribution
				% =========================================================
				\coordinate (AbusTop)    at (3.65,0 |- g1.input 1);
				\coordinate (AbusBottom) at (3.65,0 |- g3.input 1);

				\draw[wire]
				(Asplit)
				-- (Asplit |- g1.input 1)
				-- (AbusTop)
				-- (g1.input 1);

				\draw[wire]
				(AbusTop)
				-- (AbusBottom)
				-- (g3.input 1);

				\node[junction] at (AbusTop) {};

				% =========================================================
				% B input
				% =========================================================
				\coordinate (Bsplit) at (0.75,0 |- Blevel);

				\draw[wire]
				(Blabel.east)
				-- (notB.input);

				\node[junction] at (Bsplit) {};

				% =========================================================
				% B' distribution
				% =========================================================
				\coordinate (BbarBase)   at (4.35,0 |- Blevel);
				\coordinate (BbarTop)    at (4.35,0 |- g0.input 2);
				\coordinate (BbarMiddle) at (4.35,0 |- g1.input 2);

				\draw[wire]
				(notB.output)
				-- (BbarBase);

				\draw[wire]
				(BbarBase)
				-- (BbarTop)
				-- (g0.input 2);

				\draw[wire]
				(BbarMiddle)
				-- (g1.input 2);

				\node[junction] at (BbarMiddle) {};

				% =========================================================
				% B distribution
				% =========================================================
				\coordinate (BbusTop)    at (4.95,0 |- g2.input 2);
				\coordinate (BbusBottom) at (4.95,0 |- g3.input 2);

				\draw[wire]
				(Bsplit)
				-- (Bsplit |- g2.input 2)
				-- (BbusTop)
				-- (g2.input 2);

				\draw[wire]
				(BbusTop)
				-- (BbusBottom)
				-- (g3.input 2);

				\node[junction] at (BbusTop) {};

				% =========================================================
				% Enable distribution
				% =========================================================
				\coordinate (ENbusBottom) at (5.75,0 |- g3.input 3);
				\coordinate (ENbus2)      at (5.75,0 |- g2.input 3);
				\coordinate (ENbus1)      at (5.75,0 |- g1.input 3);
				\coordinate (ENbusTop)    at (5.75,0 |- g0.input 3);

				\draw[wire]
				(ENlabel.east)
				-- (ENbusBottom)
				-- (g3.input 3);

				\draw[wire]
				(ENbusBottom)
				-- (ENbusTop)
				-- (g0.input 3);

				\draw[wire]
				(ENbus1)
				-- (g1.input 3);

				\draw[wire]
				(ENbus2)
				-- (g2.input 3);

				\node[junction] at (ENbusBottom) {};
				\node[junction] at (ENbus1) {};
				\node[junction] at (ENbus2) {};

				% =========================================================
				% Outputs
				% =========================================================
				\draw[wire]
				(g0.output)
				-- ++(1.35,0)
				node[right] {$A'B'$};

				\draw[wire]
				(g1.output)
				-- ++(1.35,0)
				node[right] {$AB'$};

				\draw[wire]
				(g2.output)
				-- ++(1.35,0)
				node[right] {$A'B$};

				\draw[wire]
				(g3.output)
				-- ++(1.35,0)
				node[right] {$AB$};

			\end{tikzpicture}%
		}
	\end{minipage}
	\hfill
	\begin{minipage}[c]{0.38\textwidth}
		\centering
		\renewcommand{\arraystretch}{1.15}
		\setlength{\tabcolsep}{5pt}
		\resizebox{\linewidth}{!}{%
			\begin{tabular}{|c|c|c|c|c|c|c|}
				\hline
				$E$ & $A$ & $B$ & $D_0$ & $D_1$ & $D_2$ & $D_3$ \\
				\hline
				0   & X   & X   & 0     & 0     & 0     & 0     \\
				\hline
				1   & 0   & 0   & 1     & 0     & 0     & 0     \\
				\hline
				1   & 0   & 1   & 0     & 1     & 0     & 0     \\
				\hline
				1   & 1   & 0   & 0     & 0     & 1     & 0     \\
				\hline
				1   & 1   & 1   & 0     & 0     & 0     & 1     \\
				\hline
			\end{tabular}%
		}
	\end{minipage}

\end{figure}
\end{problem}

\begin{solution}
	Number the outputs as in the truth table: $D_0 = E A'B'$, $D_1 = E A'B$,
	$D_2 = E AB'$, $D_3 = E AB$, i.e.\ the four AND gates of the figure, each
	gated by $EN$. Fault free, $(D_0,D_1,D_2,D_3)$ carries exactly one $1$
	when $E=1$ and is all-zero when $E=0$. The first half of that is a
	$1$-out-of-$4$ code, which is easy to check, but the idle word $0000$ is
	not in it, so a checker watching the four outputs alone would report a
	fault every time the demux is disabled.

	\myspace
	\textbf{Step 1: make the output code constant weight.} Add a fifth
	``idle'' line $D_4 = E'$ (one inverter on $EN$). Then
	$(D_0,D_1,D_2,D_3,D_4)$ is a $1$-out-of-$5$ codeword for \emph{every}
	input combination, with $D_4=1$ marking the disabled state. Any fault
	that kills the active output, turns on a second one, or lets an output
	through while $E=0$ now pushes the word off the code (weight $0$ or $2$).

	\myspace
	\textbf{Step 2: reduce $1$-out-of-$5$ to two-rail signals.} A pair of OR
	gates over a subset of the lines and its complement produces a two-rail
	(complementary) pair for every codeword. Such a pair misses a
	double-active word only when both active lines sit on the same side of
	the split, so with $k$ pairs each line carries a $k$-bit membership
	signature and two lines are distinguished only if their signatures
	differ: $2^k \ge 5$, hence $k = 3$ pairs is the minimum. Splitting on
	``idle'', on $A$ and on $B$ gives three pairs, five OR gates in total:
	\[
		P = \paren{D_0{+}D_1{+}D_2{+}D_3,\ D_4}, \qquad
		Q = \paren{D_0{+}D_1{+}D_4,\ D_2{+}D_3}, \qquad
		R = \paren{D_0{+}D_2{+}D_4,\ D_1{+}D_3}.
	\]

	\myspace
	\textbf{Step 3: merge the pairs with two-rail checker cells.} A TRC cell
	takes two two-rail pairs $(a_0,a_1)$, $(b_0,b_1)$ and outputs
	\[
		c_0 = a_0b_0 + a_1b_1, \qquad c_1 = a_0b_1 + a_1b_0,
	\]
	which is complementary iff both inputs are. Cascading two cells,
	$S = \mathrm{TRC}(P,Q)$ and $Z = \mathrm{TRC}(S,R)$, gives one final
	pair: $Z_0 \neq Z_1$ means healthy, $Z_0 = Z_1$ ($00$ or $11$) means
	fault.

	\begin{figure}[H]
		\centering
		\resizebox{0.95\linewidth}{!}{%
			\begin{tikzpicture}[
					>=Stealth, thick,
					grp/.style={draw, rectangle, minimum width=5.0cm, minimum height=1.0cm, font=\small},
					trc/.style={draw, rectangle, minimum width=1.6cm, minimum height=1.7cm, font=\small},
					lbl/.style={font=\small},
				]
				\node[lbl] (bus) at (0,0) {$\begin{array}{c} D_0,\,D_1,\,D_2,\,D_3 \\[2pt] D_4 = E' \end{array}$};

				\node[grp] (P) at (5.4,2.5)  {$P = (\,D_0{+}D_1{+}D_2{+}D_3,\ D_4\,)$};
				\node[grp] (Q) at (5.4,0.0)  {$Q = (\,D_0{+}D_1{+}D_4,\ D_2{+}D_3\,)$};
				\node[grp] (R) at (5.4,-2.5) {$R = (\,D_0{+}D_2{+}D_4,\ D_1{+}D_3\,)$};

				\draw[->] (bus) -- (P);
				\draw[->] (bus) -- (Q);
				\draw[->] (bus) -- (R);

				\node[trc] (T1) at (10.4,1.25)  {TRC$_1$};
				\node[trc] (T2) at (13.6,-0.65) {TRC$_2$};

				\draw[->] (P.east) -- ++(0.9,0) |- ($(T1.west)+(0,0.55)$);
				\draw[->] (Q.east) -- ++(0.9,0) |- ($(T1.west)+(0,-0.55)$);
				\draw[->] (T1.east) -- ++(0.8,0) node[lbl, above, midway]{$S$} |- ($(T2.west)+(0,0.55)$);
				\draw[->] (R.east) -- ++(4.1,0) |- ($(T2.west)+(0,-0.55)$);
				\draw[->] (T2.east) -- ++(1.6,0) node[lbl, right]{$(Z_0,Z_1)$};
			\end{tikzpicture}%
		}
		\caption{Checker for the demultiplexer: the $1$-out-of-$5$ output word
			is compressed into three two-rail pairs and merged by two two-rail
			checker cells. $Z_0 \neq Z_1$ is fault free, $Z_0 = Z_1$ signals a
			fault.}
		\label{fig:selfcheck-demux}
	\end{figure}

	\myspace
	\textbf{Fault-free behaviour.} All five normal words map to a
	complementary $Z$, and both values of $Z$ occur:

	\begin{simpletable}{cccccccc}{Checker response over the five normal input states (pairs written $x_0x_1$).}{selfcheck-demux}
		$E$ & $A$ & $B$ & Active line & $P$ & $Q$ & $R$ & $Z$ \\
		\midrule
		0 & X & X & $D_4$ & 01 & 10 & 10 & 01 \\
		1 & 0 & 0 & $D_0$ & 10 & 10 & 10 & 10 \\
		1 & 0 & 1 & $D_1$ & 10 & 10 & 01 & 01 \\
		1 & 1 & 0 & $D_2$ & 10 & 01 & 10 & 01 \\
		1 & 1 & 1 & $D_3$ & 10 & 01 & 01 & 10 \\
	\end{simpletable}

	\myspace
	\textbf{Coverage.} Every single stuck-at fault inside the demux -- on an
	inverter, on an AND gate input or output, or on an output line -- either
	drops the active output, raises an extra one, or lets one through while
	$E=0$. In all three cases the word has weight $0$ or $2$, at least one of
	$P,Q,R$ becomes non-complementary, and the TRC cells propagate that to
	$Z_0 = Z_1$. For instance $D_3$ stuck-at-$1$ while $D_0$ is active gives
	$R = (1,1)$, and the $A$ branch stuck-at-$0$ with $E=1$ silences every
	output and gives $P = (0,0)$. The checker is itself self-testing: the
	five words above drive both rails of every OR group and both values of
	$Z$, so a single stuck-at fault inside the checker also shows up as a
	non-code $Z$ for at least one normal input -- the pair is totally
	self-checking. The one thing the output code cannot see is a fault on $A$
	or $B$ \emph{upstream} of the demux: the wrong output then fires alone,
	which is still a legal codeword, so catching that requires the select
	inputs to arrive in two-rail form and be checked separately.
\end{solution}
% ================================================================================ %

\end{document}
